\documentclass[11pt]{article}
\usepackage[margin=1in]{geometry}
\usepackage{amsmath,amssymb,amsthm}
\usepackage{graphicx}
\usepackage[hidelinks]{hyperref}

\newtheorem{theorem}{Theorem}
\newtheorem{corollary}{Corollary}
\newtheorem{lemma}{Lemma}

\newcommand{\F}{\mathcal F}
\newcommand{\N}{\mathbb N}
\newcommand{\R}{\mathbb R}
\newcommand{\diam}{\operatorname{diam}}

\title{A Diamond Obstruction for Low Fragmentability of Free Spaces over Nets}
\author{Run \texttt{fa\_banach\_001}, agent \texttt{agent\_lane\_17}}
\date{2026-06-20}

\begin{document}
\maketitle

\section*{Status}

This is a partial result, not a full solution of the open problem in
Remark~3.7 of Basset--Lancien--Prochazka~\cite{BLP}.  It proves a broad
high-index obstruction for genuine nets in Banach spaces.  It does not compute
the weak-fragmentability index of the canonical integer grid of $\ell_1$.

\section*{Source Problem}

The source paper studies the weak-fragmentability index $\Phi$ and the
dentability index of Lipschitz-free spaces over countable complete discrete
metric spaces.  On PDF page 10, Remark~3.7 asks for examples of nets in
infinite-dimensional Banach spaces, or integer grids of infinite
semi-normalized Schauder bases, for which the free space has
weak-fragmentability index $\omega$.  The remark points to the canonical
integer grid of $\ell_1$ as a natural candidate.

\begin{center}
\includegraphics[width=\linewidth]{figures/open_problem_crop.png}
\end{center}

\section*{Idea of the Proof}

The source paper already proves the two index estimates that matter here.
First, every uniformly discrete metric space $M$ satisfies
$\Phi(\F(M))\leq \omega^2$; if $M$ is infinite, only $\omega$ and $\omega^2$
can occur below this bound.  Second, if the countably branching diamond graphs
$(D_n^\omega)_{n\in\N}$ embed into $M$ with one uniform distortion bound, then
$\Phi(\F(M))\geq\omega^2$.

Thus it is enough to know when the diamond family survives inside a net.  A
net projection changes distances by an additive constant.  In a Banach space,
however, a bi-Lipschitz copy can first be dilated by an arbitrarily large
scalar.  After that dilation, nearest-net projection only perturbs distances
by a small relative amount.  Since graph distances in $D_n^\omega$ are
integer-valued and at least one between distinct vertices, one scale absorbs
the additive error for all pairs and all $n$.

\section*{Definitions}

An $(a,\rho)$-net in a Banach space $X$ means a subset $N\subset X$ such that
$\|x-y\|\geq a$ for distinct $x,y\in N$, and every point of $X$ lies within
distance $\rho$ of some point of $N$.  The countably branching diamond graphs
$D_n^\omega$ are used with their graph metrics $d_n$, normalized so that the
shortest nonzero distance is $1$.

A family $(D_n^\omega,d_n)$ embeds with equal distortion into a metric space
$M$ if there is a constant $C$ such that each $D_n^\omega$ admits a
bi-Lipschitz embedding into $M$ with distortion at most $C$.

\begin{lemma}[scaled nearest-net projection]
Let $X$ be a Banach space and let $N$ be an $(a,\rho)$-net in $X$.  Suppose
that the family $(D_n^\omega,d_n)$ embeds with equal distortion into $X$.
Then $(D_n^\omega,d_n)$ embeds with equal distortion into $N$.
\end{lemma}

\begin{proof}
After rescaling the individual embeddings, choose $L\geq 1$ and maps
$f_n:D_n^\omega\to X$ such that for every $u,v\in D_n^\omega$,
\[
        d_n(u,v)\leq \|f_n(u)-f_n(v)\|\leq Ld_n(u,v).
\]
Fix a scalar $R\geq \max\{1,4\rho\}$.  For each vertex $v\in D_n^\omega$,
choose $g_n(v)\in N$ with
\[
        \|g_n(v)-R f_n(v)\|\leq \rho .
\]
If $u\neq v$, then $d_n(u,v)\geq 1$, so
\[
 \|g_n(u)-g_n(v)\|
 \geq R\|f_n(u)-f_n(v)\|-2\rho
 \geq Rd_n(u,v)-2\rho
 \geq \frac R2 d_n(u,v).
\]
Also,
\[
 \|g_n(u)-g_n(v)\|
 \leq R\|f_n(u)-f_n(v)\|+2\rho
 \leq RLd_n(u,v)+2\rho
 \leq R\left(L+\frac12\right)d_n(u,v).
\]
Thus each $g_n$ is bi-Lipschitz, with distortion at most $2L+1$,
independently of $n$.  This is the required equal-distortion embedding into
$N$.
\end{proof}

\begin{theorem}[net diamond obstruction]
Let $X$ be a Banach space and let $N$ be an $(a,\rho)$-net in $X$.  If the
countably branching diamond graphs $(D_n^\omega)_{n\in\N}$ embed with equal
distortion into $X$, then
\[
        \Phi(\F(N))=\omega^2 .
\]
\end{theorem}

\begin{proof}
By the lemma, the family $(D_n^\omega)_{n\in\N}$ embeds with equal distortion
into $N$.  Proposition~3.2 of~\cite{BLP} then yields
$\Phi(\F(N))\geq \omega^2$.

Since $N$ is a net, it is uniformly discrete.  Proposition~3.1 of~\cite{BLP}
gives $\Phi(\F(N))\leq \omega^2$.  The two estimates coincide, so
$\Phi(\F(N))=\omega^2$.
\end{proof}

\begin{corollary}
Let $X$ be a Banach space containing a linear isomorphic copy of a Banach
space into which the family $(D_n^\omega)$ embeds with equal distortion.  Then
every net $N$ in $X$ satisfies $\Phi(\F(N))=\omega^2$.
In particular, this applies to Banach spaces containing an isomorphic copy of
$c_0$, and to spaces containing any of the $L_1$ diamond models cited in
the source paper.
\end{corollary}

\begin{proof}
Compose the uniform diamond embeddings with the linear isomorphic inclusion
into $X$.  This changes the common distortion bound by at most the Banach
space isomorphism constant.  The theorem applies.
\end{proof}

\section*{Limitations}

This result does not decide whether any integer grid of an infinite
semi-normalized Schauder basis has free-space weak-fragmentability index
$\omega$.  In particular, it does not compute $\Phi(\F(\mathbb Z_{\ell_1}))$.
The proof uses coarse density of a genuine net in an ambient Banach space; an
integer grid such as the canonical grid of $\ell_1$ need not be a net in this
sense.

The result should therefore be read as a large-family exclusion theorem:
whenever the ambient Banach space already contains uniformly distorted
countably branching diamonds, no genuine net in that space can provide the
low-index example sought in Remark~3.7.

\section*{Verification and Novelty Check}

The proof is analytic and uses no finite search.  The evidence crop was
generated by
\[
\texttt{conda run --no-capture-output -n sandbox python code/make\_open\_problem\_crop.py}.
\]

Bounded duplicate and novelty checks were performed on 2026-06-20.  The local
run indexes were searched for arXiv:2505.19891, weak-fragmentability index,
countably branching diamonds, the integer grid of $\ell_1$, and nearby net
keywords.  Web searches for the source arXiv id and for combinations of
\emph{weak-fragmentability index}, \emph{Lipschitz-free space},
\emph{integer grid}, $\ell_1$, and \emph{countably branching diamond} found
the source paper but no later computation of the $\ell_1$ grid and no duplicate
of this abstract net-obstruction packet.

\section*{Human Review Recommendation}

Review as a likely valid partial result.  The central point to check is the
scaled nearest-net projection lemma, especially the use of the graph-metric
normalization $d_n(u,v)\geq 1$ for distinct vertices.  If accepted, the packet
should be treated as an obstruction/family-classification result, not as a
solution of the source paper's proposed $\mathbb Z_{\ell_1}$ candidate.

\begin{thebibliography}{9}

\bibitem{BLP}
E.~Basset, G.~Lancien, and A.~Prochazka,
\emph{Diversity of Lipschitz-free spaces over countable complete discrete
metric spaces}, arXiv:2505.19891, 2025.

\end{thebibliography}

\end{document}
